\documentclass[a4paper,12pt,oneside]{article}

\usepackage[T2A]{fontenc}
\usepackage[utf8]{inputenc}
\usepackage[russian]{babel}

\usepackage{tempora}
\usepackage{microtype}

\usepackage{amsmath,amssymb,amsthm,mathtools}

\usepackage{geometry}
\geometry{top=25mm,left=25mm,right=20mm,bottom=25mm}

\usepackage{enumitem}
\setlist[enumerate]{leftmargin=*, itemsep=0.75\baselineskip, topsep=0.35\baselineskip}
\setlist[itemize]{leftmargin=*, itemsep=0.25\baselineskip, topsep=0.25\baselineskip}

\setlength{\parindent}{1.25em}
\setlength{\parskip}{0pt}

\newcommand{\Dom}{\mathrm{Dom}}
\newcommand{\Im}{\mathrm{Im}}

\newcommand{\DefItem}[2]{\item \textbf{#1}\\#2}

\begin{document}

\begin{center}
    {\Large\bfseries Теоретический минимум}\\[2mm]
    {\normalsize 1 модуль}
\end{center}
\vspace{2mm}
\hrule
\vspace{4mm}

\begin{enumerate}[label=\arabic*.]

\DefItem{Множество.}{Множество — совокупность объектов произвольной природы, рассматриваемая как единое целое; объекты называются элементами множества.
\begin{itemize}
\item \textbf{Определение 1:} набор (совокупность) объектов, объединённых общим свойством.
\item \textbf{Определение 2:} совокупность объектов произвольной природы, объединённых общим свойством.
\item \textbf{Определение 3:} неупорядоченный набор уникальных элементов.
\end{itemize}}

\DefItem{Способы задания множеств:}{
\begin{itemize}
\item перечисление;
\item характеристическое свойство (правило);
\item порождающая процедура;
\item формула (аналитически, через операции);
\item диаграмма Венна.
\end{itemize}}

\DefItem{Подмножество и отличие собственных и несобственных подмножеств.}{$A$ — подмножество $B$, если каждый элемент $A$ принадлежит $B$. Обозначается как $A \subseteq B$.
\begin{itemize}
\item \textbf{Несобственные:} само множество $A$ и пустое $\varnothing$.
\item \textbf{Собственные:} любые подмножества множества $A$, отличные от $A$ и пустого $\varnothing$.
\end{itemize}}

\DefItem{Булеан множества.}{Булеан множества $A$ — $P(A)$ или $2^A$, множество всех подмножеств $A$.}

\DefItem{Объединение множеств.}{$A \cup B$ — множество элементов, принадлежащих хотя бы одному из множеств $A$ или $B$:
\[
A \cup B = \{x \mid x \in A \ \text{или}\ x \in B\}.
\]}

\DefItem{Пересечение множеств.}{$A \cap B$ — множество элементов, принадлежащих одновременно и $A$, и $B$:
\[
A \cap B = \{x \mid x \in A \ \text{и}\ x \in B\}.
\]}

\DefItem{Дополнение множества.}{$\overline{A}$ — дополнение множества $A$ до универсума $U$:
\[
\overline{A} = U \setminus A = \{x \mid x \in U \ \text{и}\ x \notin A\}.
\]}

\DefItem{Разность множеств.}{$A \setminus B$ — множество элементов, принадлежащих $A$ и не принадлежащих $B$:
\[
A \setminus B = \{x \mid x \in A \ \text{и}\ x \notin B\}.
\]}

\DefItem{Симметрическая разность множеств.}{$A \Delta B$ — множество элементов, принадлежащих либо только $A$, либо только $B$:
\[
A \Delta B = \{x \mid (x \in A \ \text{и}\ x \notin B)\ \text{или}\ (x \in B \ \text{и}\ x \notin A)\}.
\]}

\DefItem{Декартово произведение множеств $A$ и $B$.}{Декартово произведение — множество всех упорядоченных пар:
\[
A \times B = \{(a,b) \mid a \in A \ \text{и}\ b \in B\}.
\]}

\DefItem{Прямое произведение четырёх множеств.}{Это множество всех кортежей длины $4$, где каждая компонента из своего множества:
\[
X_1 \times X_2 \times X_3 \times X_4
=
\{(x_1,x_2,x_3,x_4)\mid x_1 \in X_1,\ x_2 \in X_2,\ x_3 \in X_3,\ x_4 \in X_4\}.
\]}

\DefItem{Кортеж (вектор).}{Упорядоченный набор произвольных объектов.}

\DefItem{Равенство упорядоченных пар/кортежей.}{Кортежи одной длины равны поэлементно; аналогично для упорядоченных пар.}

\DefItem{Свойство коммутативности.}{Операции объединения и пересечения коммутативны:
\[
A \cup B = B \cup A, \qquad A \cap B = B \cap A.
\]}

\DefItem{Свойство дистрибутивности.}{Верны тождества:
\[
(A \cap B) \cup C = (A \cup C)\cap(B \cup C),
\qquad
(A \cup B) \cap C = (A \cap C)\cup(B \cap C).
\]}

\DefItem{Критерий равенства двух множеств (метод двойного включения).}{
\[
A=B \iff (A \subseteq B)\ \text{и}\ (B \subseteq A).
\]}

\DefItem{Равные множества.}{$A = B$ тогда и только тогда, когда состоят из одних и тех же элементов:
\[
\forall x:\ x \in A \iff x \in B.
\]}

\DefItem{Свойство ассоциативности.}{Операции объединения и пересечения ассоциативны:
\[
A \cup (B \cup C) = (A \cup B) \cup C,
\qquad
A \cap (B \cap C) = (A \cap B) \cap C.
\]}

\DefItem{Закон поглощения.}{
\[
(A \cap B)\cup A = A,
\qquad
(A \cup B)\cap A = A.
\]}

\DefItem{Закон склеивания.}{
\[
(A \cap B)\cup(A \cap \overline{B})=A,
\qquad
(A \cup B)\cap(A \cup \overline{B})=A.
\]}

\DefItem{Законы де Моргана (множества).}{
\[
\overline{A \cup B} = \overline{A} \cap \overline{B},
\qquad
\overline{A \cap B} = \overline{A} \cup \overline{B}.
\]}

\DefItem{Бинарное отношение на множествах.}{Бинарное отношение $R$ на паре множеств $A$ и $B$ — подмножество $A \times B$. Если $(a,b)\in R$, то пишут $aRb$. Если $A=B$, то $R\subseteq A\times A$ — бинарное отношение на $A$. Всего возможных отношений:
\[
2^{|A\times B|}.
\]}

\DefItem{Способы задания бинарных отношений.}{
\begin{itemize}
\item правило;
\item перечисление;
\item таблица;
\item на плоскости (система координат);
\item матрица;
\item граф;
\item графический способ.
\end{itemize}}

\DefItem{Область определения и область значений.}{Для $R\subseteq A\times B$:
\[
\Dom(R)=\{x \in A \mid \exists\,y \in B:\ (x,y)\in R\},
\qquad
\Im(R)=\{y \in B \mid \exists\,x \in A:\ (x,y)\in R\}.
\]}

\DefItem{Обратное отношение.}{Обратное отношение $R^{-1}$ — отношение на $B \times A$:
\[
R^{-1}=\{(b,a)\mid (a,b)\in R\}.
\]}

\DefItem{Композиция отношений.}{Пусть $R\subseteq A\times B$ и $S\subseteq B\times C$. Тогда композиция — отношение $T$ на $A\times C$:
\[
T=\{(a,c)\mid \exists\,b \in B:\ (a,b)\in R \ \text{и}\ (b,c)\in S\}.
\]
Обозначают: $T=S\circ R$ (также встречается $SR$, $R;S$).}

\DefItem{Степень отношения.}{Для $R\subseteq A\times A$:
\[
R^n = R^{n-1}\circ R,\qquad n\ge 2.
\]}

\DefItem{Теорема об ассоциативности композиции отношений.}{Если $R\subseteq A\times B$, $S\subseteq B\times C$, $T\subseteq C\times D$, то
\[
(T\circ S)\circ R = T\circ (S\circ R).
\]}

\DefItem{Симметричность.}{$\forall a,b\in A:\ aRb \Rightarrow bRa$.}

\DefItem{Асимметричность.}{$\forall a,b\in A:\ aRb \Rightarrow \neg(bRa)$.}

\DefItem{Антисимметричность.}{$\forall a,b\in A:\ (aRb\ \text{и}\ bRa)\Rightarrow a=b$.}

\DefItem{Несимметричность.}{Отношение называется несимметричным, если оно не является симметричным и не является асимметричным.}

\DefItem{Рефлексивность.}{$\forall a\in A:\ aRa$.}

\DefItem{Антирефлексивность.}{$\forall a\in A:\ \neg(aRa)$.}

\DefItem{Транзитивность.}{$\forall a,b,c\in A:\ (aRb\ \text{и}\ bRc)\Rightarrow aRc$.}

\DefItem{Нетранзитивность.}{Существуют $a,b,c\in A$ такие, что $aRb$ и $bRc$, но $\neg(aRc)$.}

\DefItem{Нерефлексивность.}{Для некоторых $a\in A$ выполняется $aRa$, а для некоторых — нет.}

\DefItem{Антитранзитивность.}{$\forall a,b,c\in A:\ (aRb\ \text{и}\ bRc)\Rightarrow \neg(aRc)$.}

\DefItem{Замыкание отношения $R$ относительно свойства $P$.}{Замыкание — минимальное по включению надмножество $R$, обладающее свойством $P$.}

\DefItem{Отношение эквивалентности $R$.}{Отношение эквивалентности — отношение, которое рефлексивно, симметрично и транзитивно.}

\DefItem{Класс эквивалентности.}{Подмножества, образованные в результате разбиения множества отношением эквивалентности.}

\DefItem{Порождающий элемент класса эквивалентности.}{Элемент, находящийся в отношении $R$ со всеми элементами соответствующего класса. Обозначения:
\[
[a]\ \text{— класс эквивалентности, порождённый элементом } a,
\qquad
[A]_R\ \text{— множество всех классов эквивалентности по отношению } R.
\]}

\DefItem{Разбиение множества.}{Пусть $\langle A\rangle=\{A_i:\ i\in I\}$ — множество непустых подмножеств $A$ таких, что
\[
A_i\cap A_j=\varnothing \ \text{для всех } i\ne j,
\qquad
A=\bigcup_{i\in I} A_i.
\]}

\DefItem{Теорема про разбиение и классы эквивалентности.}{Непустое множество подмножеств $\langle A\rangle$ есть разбиение $A$ тогда и только тогда, когда $\langle A\rangle=[A]_R$ по некоторому отношению эквивалентности $R$.}

\DefItem{Отношение строгого и нестрогого частичного порядка.}{
\begin{itemize}
\item \textbf{Строгий:} транзитивность, асимметричность, антирефлексивность.
\item \textbf{Нестрогий:} транзитивность, антисимметричность, рефлексивность.
\end{itemize}}

\DefItem{Отношение строгого и нестрогого линейного порядка.}{
\begin{itemize}
\item \textbf{Строгий:} для любых двух $a\ne b$ существует либо $aRb$, либо $bRa$.
\item \textbf{Нестрогий:} для любых двух $a,b$ существует либо $aRb$, либо $bRa$.
\end{itemize}}

\DefItem{Отношение соответствия.}{
\begin{itemize}
\item \textbf{Одно-однозначное соответствие.} Каждому $a\in A$ соответствует ровно один $b\in B$, и каждому $b\in B$ соответствует ровно один $a\in A$.
\item \textbf{Одно-многозначное соответствие.} Каждому $a\in A$ соответствует несколько $b\in B$, но каждому $b\in B$ соответствует только один $a\in A$.
\item \textbf{Много-однозначное соответствие.} Каждому $a\in A$ соответствует ровно один $b\in B$, но некоторому $b\in B$ соответствуют несколько $a\in A$.
\item \textbf{Много-многозначное соответствие.} Каждому $a\in A$ соответствует несколько $b\in B$, и каждому $b\in B$ соответствует несколько $a\in A$.
\end{itemize}}

\DefItem{Линейно упорядоченное множество.}{Линейно упорядоченное множество — множество, на котором задано отношение линейного порядка: $\langle M,R\rangle$.}

\DefItem{Частично упорядоченное множество.}{Частично упорядоченное множество — множество, упорядоченное строгим или нестрогим порядком, в котором могут существовать несравнимые элементы (не для всех пар $a,b$ верно $aRb$ или $bRa$).}

\DefItem{Функциональное отношение.}{Функциональное отношение — бинарное отношение $xRy$, в котором каждому $x\in X$ соответствует не более одного $y\in Y$; запись: $y=F(x)$, где $F \subset X \times Y$.}

\DefItem{Недоопределенная функция.}{Если функция определена не на всём множестве $X$, то такая функция будет недоопределённой, то есть $\Dom(f)\ne X$.}

\DefItem{Какими свойствами обладает отношение, когда нельзя определить симметричность однозначно?}{Симметрично и антисимметрично, или антисимметрично и асимметрично.}

\DefItem{Какими свойствами обладает отношение, когда нельзя определить транзитивность однозначно?}{Транзитивно и антитранзитивно.}

\end{enumerate}

\end{document}
